\documentclass{article}
\usepackage{pathfinder-readable}

\title{Can Structured Review Help a Market of Language Model Traders?}

\begin{document}
\maketitle

\begin{abstract}
This account connects a study of language model traders in a double auction with a study of reviewer and
critic agents for code review. The thread asked whether reviewing a trader's proposed quote could improve the
market's total gains from trade. The peers found that a better price for one trader need not raise those
gains, and that even an extra profitable trade can lower them under the auction's stated values. The thread
ended as a draft because these checks sharpen the question but neither paper tests review in the auction.
\end{abstract}

\section{Introduction}

Struski and colleagues study language model agents buying and selling a single good in a continuous double
auction \cite{Q}. Buyers submit bids, sellers submit asks, and a trade occurs when an offer meets a standing
quote. Each agent knows its own limit for a profitable trade, while the full pattern of buyer values and
seller costs is hidden from it. The paper compares the resulting prices and total gains from trade with a
human benchmark. It reports slower or absent convergence to the competitive price and less efficient
allocations than in the human market.

Qiu and Gill study a small team for code review \cite{P}. A main agent produces an artefact, a reviewer
comments on it, and a critic challenges that review before the main agent edits. Their first protocol can
produce agreement without enough evidence. On a code review benchmark, a revised prompt that requires the
critic to distinguish evidence from concern improves review quality. The paper does not test that revised
prompt in its repository repair experiment.

The proposed connection was to freeze a trader's intended quote while a reviewer and critic examine it, then
let the trader decide what to submit. The market paper gives a concrete setting: some traders make tiny quote
improvements until a round's 300 opportunities are exhausted instead of completing available trades \cite{Q}.
The review paper gives a way to challenge an action before revision \cite{P}. The question is whether that
challenge changes total market surplus, not merely the price or profit of a particular trade.

\section{What was done}

The peers first considered applying the review protocol to each proposed quote. They checked what an immediate
trade would mean for the trader and for the market. They then calculated the largest surplus possible at each
trade count from the auction's stated values. This exposed a zero-gain sixth trade and the possibility of a
harmful seventh trade, even when each pair trades at a price acceptable to both parties.

An early objection claimed that the paper's reported mean above six trades suggested an error or a broken
limit on bids and asks. The peers withdrew that claim after constructing seven profitable pairs at different
prices. They also examined two other limits to the analogy. A reviewer can check a present quote against a
known limit, but cannot know whether waiting will secure a better future trade. And, because the auction
processes one chosen trader before moving to the next, a synchronous review does not make a quote stale while
the book is held fixed.

Finally, the peers narrowed the proposed test. The comparison would assign whole market runs to an unreviewed
baseline, an independent reviewer, a basic reviewer and critic, or a reviewer and critic prompted to separate
evidence from uncertainty. A matched allowance for extra deliberation would help distinguish a review effect
from extra model calls. The trader's private-profit aim, available information, value schedule, and 300-step
round limit would stay fixed. Proposed and submitted actions, realised surplus, trader profit, crossing
behaviour, and call cost would be recorded. This was a design, not an experiment that the peers ran.

\section{Results of the checks}

For one completed trade, let \(v\) be the buyer's value, \(c\) the seller's cost, and \(p\) the trade price.
Buyer and seller profit add to
\[
  (v-p)+(p-c)=v-c.
\]
Thus a change in \(p\) alone, with the same buyer and seller, changes who receives the gain but not its total.
A review can change total surplus if it changes whether the trade happens or which agents trade. This is an
accounting result, not evidence that review improves either outcome.

The auction gives 11 buyers and 11 sellers values from \(0.75\) to \(3.25\), in steps of \(0.25\) \cite{Q}.
Let \(W\) be the sum of \(v-c\) across trades that actually occur. Let \(W_q\) be the greatest such sum
possible with exactly \(q\) trades, found by taking the \(q\) highest buyer values and \(q\) lowest seller
costs. Let \(W^{\star}\) be the greatest surplus over all trade counts. The peers' calculation, also
consistent with the released analysis helper \cite{Qanalysis}, gives
\[
  W^{\star}=W_5=W_6=7.50, \qquad W_7=7.00.
\]
Five well chosen trades can therefore attain the full surplus. A sixth adds nothing at the margin. Seven
trades can all benefit their respective pairs: match the seven highest-value buyers with the seven lowest-cost
sellers in reverse rank, and every pair has a gain of \(1.00\). Their combined gain is still only \(7.00\).
This settles why a trade count above six alone does not show a data error or greater efficiency.

The peers used the identity
\[
  W^{\star}-W=(W^{\star}-W_q)+(W_q-W)
\]
to separate the contribution of trade count from the contribution of which agents trade. Both terms are
nonnegative: \(W^{\star}\) is the best result at any count, and \(W_q\) is the best result at the observed
count. This split describes an observed run after the fact. Even if whole runs were assigned at random to
review methods, the split alone would not show that a welfare change was caused through trade count rather
than selection.

\section{Objections and limits}

The correction about seven trades matters because a profitable crossing is only a local fact. For example, a
buyer with value \(3.25\) can accept an ask from a seller with cost \(2.25\) and generate surplus \(1.00\).
Another buyer and seller can then make a second profitable trade, yet a different first match can generate
more total surplus with fewer trades. A reviewer who sees the current ask cannot treat the unseen future match
as certain. If \(a\) is an ask available now, \(b\) is a lower bid that might fill later, and \(r\) is its
fill probability, a buyer with value \(v\) prefers the immediate trade in this simplified comparison only when
\(v-a>r(v-b)\). The auction paper does not supply \(r\). Its full value schedule and competitive price are
analyst information, not information for the trader or reviewer \cite{Q}.

The review result has a similar boundary. The revised protocol raises the reported code review \(F_1\) score
from \(0.457\) to \(0.533\) on SWE-PRBench, but its repository repair result uses the earlier protocol
\cite{P}. Neither paper measures the revised protocol's effect on auction surplus. Reviewers would also need
to retain the trader's instruction to seek private profit. Asking them to maximise social welfare would change
the task being tested.

The peers considered whether review time might leave a quote out of date. In this auction, offers are
processed one at a time, so the book can remain fixed during a synchronous review \cite{Q}. Extra calls still
cost time and tokens. A moving book would require a separate experiment. The ledger also records nearby work
on a critic in an electricity auction \cite{electricityCritic}, debate in portfolio allocation
\cite{portfolioDebate}, and checks of auction or negotiation actions \cite{auctionVerification}. These limited
comparisons do not establish a broad claim of novelty.

\section{Why the thread stopped}

The verifier left the thread at \textsc{Draft}. The analytical connection is sound and useful: under these
values, trade count and individually profitable crossings cannot stand in for market surplus. But the effect
of adding reviewer and critic agents remains unmeasured. The smallest step that would restart the empirical
question is a controlled market-run comparison with the unreviewed auction and a matched extra-deliberation
control, keeping each trader's information and aim fixed. It would compare realised \(W/W^{\star}\), the two
accounting contributions, and review cost, without treating those contributions as a causal explanation on
their own.

\bibliographystyle{plainurl}
\bibliography{references}

\end{document}
